% GENERATED FILE. Edit canonical lessons, then run npm run generate:books.
% canonical-source-hash: fnv1a64:e3da16bc16fd6ff8
% canonical-lessons: JA-C453-shiryou, JA-C453-uchiawase, JA-C453-shimekiri, JA-C453-teishutsu-suru, JA-C453-tetsuya-suru

\chapter{Handouts and Deadlines}
\label{ch:ja-handouts-and-deadlines}
\hlchaptermodality{453}

\begin{chapteropening}
\textbf{By the end of this chapter:} \emph{I can talk about handouts, meetings and deadlines.}

Complete the last lesson of chapter 453: I can talk about handouts, meetings and deadlines.
\end{chapteropening}

\section[\texorpdfstring{\ja{しりょう} (shiryō)}{しりょう (shiryō)}]{\ja{しりょう} (shiryō) — handouts; materials}

\label{lesson:JA-C453-shiryou}



\begin{quote}
Before the new one: say the Japanese for a company employee, then the Japanese for a business card.
\end{quote}

\subsection*{You'll want to know: \ja{しりょう}}
\textbf{\ja{しりょう}} — \emph{shiryō} — \textquotedblleft{}handouts; materials\textquotedblright{}.

\begin{quote}
\textbf{\ja{かいぎの} \ja{しりょうを} \ja{みんなに} \ja{わたしました。}} — \emph{kaigi no shiryō o minna ni watashimashita} — \textquotedblleft{}I gave everyone the handouts for the meeting.\textquotedblright{}
\end{quote}

\begin{cousinweb}
\textbf{\ja{しりょう}} (\emph{shiryō}) is Sino-Japanese (kango): two Chinese characters read with their on'yomi, readings borrowed from Chinese. They are \emph{shi}, \textquotedblleft{}resources\textquotedblright{} (Middle Chinese \emph{tsij}) and \emph{ryō}, \textquotedblleft{}materials\textquotedblright{} (\emph{lew}) (KANJIDIC2; Middle Chinese in Baxter's transcription).
\end{cousinweb}

\begin{grammarlens}[title={What you've built}]
One more word for describing what you lived and felt.
\end{grammarlens}

\subsection*{Your turn}
\begin{itemize}
\raggedright
  \item \emph{Say it:} \emph{shiryō}
  \item \emph{Say it:} \emph{shiryō}, once more
  \item \emph{Say it:} \emph{kaigi no shiryō o minna ni watashimashita} — I gave everyone the handouts for the meeting.
  \item \emph{Recall:} say the Japanese for a company employee, then the Japanese for a business card, then say \emph{shiryō} again
\end{itemize}

\begin{tcolorbox}[breakable,colback=teal!4,colframe=teal!35!black,title={Before you move on}]
What does \ja{しりょう} mean? (\textquotedblleft{}Handouts; materials\textquotedblright{}.) Say it once more.
\end{tcolorbox}
\section[\texorpdfstring{\ja{うちあわせ} (uchiawase)}{うちあわせ (uchiawase)}]{\ja{うちあわせ} (uchiawase) — a briefing; a planning meeting}

\label{lesson:JA-C453-uchiawase}



\begin{quote}
Before the new one: say the Japanese for a business card, then the Japanese for handouts.
\end{quote}

\subsection*{You'll want to know: \ja{うちあわせ}}
\textbf{\ja{うちあわせ}} — \emph{uchiawase} — \textquotedblleft{}a briefing; a planning meeting\textquotedblright{}.

\begin{quote}
\textbf{\ja{あしたの} \ja{あさ、うちあわせが} \ja{あります。}} — \emph{ashita no asa, uchiawase ga arimasu} — \textquotedblleft{}There is a briefing tomorrow morning.\textquotedblright{}
\end{quote}

\begin{cousinweb}
\textbf{\ja{うちあわせ}} (\emph{uchiawase}) is a native Japanese compound (wago), inherited rather than borrowed from Chinese: \emph{uchi}, \textquotedblleft{}striking (utsu)\textquotedblright{} and \emph{awase}, \textquotedblleft{}fitting together (awaseru)\textquotedblright{} (JMdict, KANJIDIC2).
\end{cousinweb}

\begin{grammarlens}[title={What you've built}]
One more word for describing what you lived and felt.
\end{grammarlens}

\subsection*{Your turn}
\begin{itemize}
\raggedright
  \item \emph{Say it:} \emph{uchiawase}
  \item \emph{Say it:} \emph{uchiawase}, once more
  \item \emph{Say it:} \emph{ashita no asa, uchiawase ga arimasu} — There is a briefing tomorrow morning.
  \item \emph{Recall:} say the Japanese for a business card, then the Japanese for handouts, then say \emph{uchiawase} again
\end{itemize}

\begin{tcolorbox}[breakable,colback=teal!4,colframe=teal!35!black,title={Before you move on}]
What does \ja{うちあわせ} mean? (\textquotedblleft{}A briefing; a planning meeting\textquotedblright{}.) Say it once more.
\end{tcolorbox}
\section[\texorpdfstring{\ja{しめきり} (shimekiri)}{しめきり (shimekiri)}]{\ja{しめきり} (shimekiri) — a cut-off date}

\label{lesson:JA-C453-shimekiri}



\begin{quote}
Before the new one: say the Japanese for handouts, then the Japanese for a briefing.
\end{quote}

\subsection*{You'll want to know: \ja{しめきり}}
\textbf{\ja{しめきり}} — \emph{shimekiri} — \textquotedblleft{}a cut-off date\textquotedblright{}.

\begin{quote}
\textbf{\ja{しめきりは} \ja{らいしゅうの} \ja{きんようびです。}} — \emph{shimekiri wa raishū no kin'yōbi desu} — \textquotedblleft{}The cut-off date is next Friday.\textquotedblright{}
\end{quote}

\begin{cousinweb}
\textbf{\ja{しめきり}} (\emph{shimekiri}) is a native Japanese compound (wago), inherited rather than borrowed from Chinese: \emph{shime}, \textquotedblleft{}fastening (shimeru)\textquotedblright{} and \emph{kiri}, \textquotedblleft{}cutting (kiru)\textquotedblright{} (JMdict, KANJIDIC2).
\end{cousinweb}

\begin{grammarlens}[title={What you've built}]
One more word for describing what you lived and felt.
\end{grammarlens}

\subsection*{Your turn}
\begin{itemize}
\raggedright
  \item \emph{Say it:} \emph{shimekiri}
  \item \emph{Say it:} \emph{shimekiri}, once more
  \item \emph{Say it:} \emph{shimekiri wa raishū no kin'yōbi desu} — The cut-off date is next Friday.
  \item \emph{Recall:} say the Japanese for handouts, then the Japanese for a briefing, then say \emph{shimekiri} again
\end{itemize}

\begin{tcolorbox}[breakable,colback=teal!4,colframe=teal!35!black,title={Before you move on}]
What does \ja{しめきり} mean? (\textquotedblleft{}A cut-off date\textquotedblright{}.) Say it once more.
\end{tcolorbox}
\section[\texorpdfstring{\ja{ていしゅつする} (teishutsu suru)}{ていしゅつする (teishutsu suru)}]{\ja{ていしゅつする} (teishutsu suru) — to hand in; to submit}

\label{lesson:JA-C453-teishutsu-suru}



\begin{quote}
Before the new one: say the Japanese for a briefing, then the Japanese for a cut-off date.
\end{quote}

\subsection*{You'll want to know: \ja{ていしゅつする}}
\textbf{\ja{ていしゅつする}} — \emph{teishutsu suru} — \textquotedblleft{}to hand in; to submit\textquotedblright{}.

\begin{quote}
\textbf{\ja{しめきりの} \ja{まえに} \ja{ていしゅつして} \ja{ください。}} — \emph{shimekiri no mae ni teishutsu shite kudasai} — \textquotedblleft{}Please hand it in before the cut-off date.\textquotedblright{}
\end{quote}

Polite \textbf{\ja{ていしゅつします}}, \emph{teishutsushimasu}; a \ja{する} verb. Past (polite): \textbf{\ja{がくせいが} \ja{しゅくだいを} \ja{ていしゅつしました。}}, \emph{gakusei ga shukudai o teishutsu shimashita}.

\begin{cousinweb}
\textbf{\ja{ていしゅつする}} (\emph{teishutsu suru}) is Sino-Japanese (kango) \emph{teishutsu} with the native verb \emph{suru}, \textquotedblleft{}to do\textquotedblright{}: two Chinese characters read with their on'yomi, readings borrowed from Chinese. They are \emph{tei}, \textquotedblleft{}carry in hand\textquotedblright{} (Middle Chinese \emph{dej}) and \emph{shutsu}, \textquotedblleft{}go out\textquotedblright{} (\emph{tsyhwit}) (KANJIDIC2; Middle Chinese in Baxter's transcription).
\end{cousinweb}

\begin{grammarlens}[title={What you've built}]
One more word for describing what you lived and felt.
\end{grammarlens}

\subsection*{Your turn}
\begin{itemize}
\raggedright
  \item \emph{Say it:} \emph{teishutsu suru}
  \item \emph{Say it:} \emph{teishutsu suru}, once more
  \item \emph{Say it:} \emph{shimekiri no mae ni teishutsu shite kudasai} — Please hand it in before the cut-off date.
  \item \emph{Recall:} say the Japanese for a briefing, then the Japanese for a cut-off date, then say \emph{teishutsu suru} again
\end{itemize}

\begin{tcolorbox}[breakable,colback=teal!4,colframe=teal!35!black,title={Before you move on}]
What does \ja{ていしゅつする} mean? (\textquotedblleft{}To hand in; to submit\textquotedblright{}.) Say it once more.
\end{tcolorbox}
\section[\texorpdfstring{\ja{てつやする} (tetsuya suru)}{てつやする (tetsuya suru)}]{\ja{てつやする} (tetsuya suru) — to stay up all night}

\label{lesson:JA-C453-tetsuya-suru}



\begin{quote}
Before the new one: say the Japanese for a cut-off date, then the Japanese for to hand in.
\end{quote}

\subsection*{You'll want to know: \ja{てつやする}}
\textbf{\ja{てつやする}} — \emph{tetsuya suru} — \textquotedblleft{}to stay up all night\textquotedblright{}.

\begin{quote}
\textbf{\ja{しけんの} \ja{まえに} \ja{てつやしました。}} — \emph{shiken no mae ni tetsuya shimashita} — \textquotedblleft{}I stayed up all night before the exam.\textquotedblright{}
\end{quote}

Polite \textbf{\ja{てつやします}}, \emph{tetsuyashimasu}; a \ja{する} verb. Past (polite): \textbf{\ja{あにが} \ja{てつやしました。}}, \emph{ani ga tetsuya shimashita}.

\begin{cousinweb}
\textbf{\ja{てつやする}} (\emph{tetsuya suru}) is Sino-Japanese (kango) \emph{tetsuya} with the native verb \emph{suru}, \textquotedblleft{}to do\textquotedblright{}: two Chinese characters read with their on'yomi, readings borrowed from Chinese. They are \emph{tetsu}, \textquotedblleft{}sit up (all night)\textquotedblright{} (Middle Chinese \emph{drjet}) and \emph{ya}, \textquotedblleft{}night\textquotedblright{} (\emph{yaeH}) (KANJIDIC2; Middle Chinese in Baxter's transcription).
\end{cousinweb}

\begin{grammarlens}[title={What you've built}]
One more word for describing what you lived and felt.
\end{grammarlens}

\subsection*{Your turn}
\begin{itemize}
\raggedright
  \item \emph{Say it:} \emph{tetsuya suru}
  \item \emph{Say it:} \emph{tetsuya suru}, once more
  \item \emph{Say it:} \emph{shiken no mae ni tetsuya shimashita} — I stayed up all night before the exam.
  \item \emph{Recall:} say the Japanese for a cut-off date, then the Japanese for to hand in, then say \emph{tetsuya suru} again
\end{itemize}

\begin{tcolorbox}[breakable,colback=teal!4,colframe=teal!35!black,title={Before you move on}]
What does \ja{てつやする} mean? (\textquotedblleft{}To stay up all night\textquotedblright{}.) Say it once more.
\end{tcolorbox}
